\section{Related Work}

Manipulation of 1D deformable objects has applications in surgical suturing~\cite{schulman_case_study_2013,suturing_autolab_2016}, knot-tying~\cite{van_den_berg_2010,case_study_knots_1991,knot_planning_2003,tying_precisely_2016}, untangling ropes~\cite{rope_untangling_2013,grannen2020untangling,tactile_cable_2020}, deforming wires~\cite{wire_insertion_1996,wire_insertion_1997}, and inserting cables into obstacles~\cite{tight_tolerance_insertion_2015}; Sanchez et al.~\cite{survey_2018} provide a survey.

Various approaches have been proposed for quasistatic manipulation of cables, such as using demonstrations~\cite{schulman_isrr_2013,lee_LFD_IROS_2014,seita_bags_2021}, self-supervised learning~\cite{nair_rope_2017,ZSVI_2018, jin2024multi}, model-free reinforcement learning~\cite{lerrel_2020,corl2020softgym, elmquist2022art}, contrastive learning~\cite{yan2020learning, sim2019personalization}, dense object descriptors~\cite{priya_2020, eisner2022flowbot3d, pan2023tax}, generative modeling~\cite{wang_visual_planning_2019, avigal20206, avigal2021avplug}, and state estimation~\cite{state_estimation_LDO, zhang2020dex, zhang2021robots, zhang2023flowbot++, yao2023apla}. These prior works generally assume quasistatic dynamics to enable a series of small deformations, often parameterized with pick-and-place actions. Alternative approaches involve sliding cables~\cite{one_hand_knotting_tactile_2007,zhu_sliding_cables_2019} or assuming non-quasistatic systems that enable tactile feedback~\cite{tactile_cable_2020, jin2024multi}. This paper focuses on dynamic actions to manipulate cables from one configuration to another, traversing over obstacles and potentially knocking over objects.

More closely related is the work of Yamakawa et al.~\cite{dynamic_knotting_2010,yamakawa2012simple,high_speed_knotting_2013,one_hand_knotting_tactile_2007, zhang2016health, devgon2020orienting, lim2021planar, lim2022real2sim2real} that focuses on dynamic manipulation of cables using robot arms moving at high speeds, enabling them to model rope deformation algebraically under the assumption that the effect of gravity is negligible. They show impressive results with a single multi-fingered robot arm that ties knots by initially whipping cables upwards. Because we use fixed-endpoint cables moving at slower speeds, the forces exerted on the rope by gravity and the fixed endpoint prevent us from performing knot ties. %assuming a model where each rope joint follows the robot motion with constant time-delay, as in Yamakawa~et~al. Instead, we use a learning-based method and do not use specialized end-effectors.
Kim and Shell~\cite{kim2016using, shen2024diffclip} use a small-scale mobile robot with a cable attached as a tail to manipulate small objects by leveraging inter-object friction to drag or strike the object of interest. While their work aims to tackle dynamic cable control, it relies heavily on physics models for motion primitives that are only applicable to dragging objects via cables, and uses a RRT-based graph search algorithm to address the stochasticity in the system. Other studies have developed physics models to analyze dynamic cable motion. For example, Zimmermann et al.~\cite{zimmermanndynamic} use the finite-element method to model elastic objects for dynamic manipulation of free-end deformable objects, and Gatti-Bonoa and Perkins~\cite{fly_fishing_2004} and Wang and Wereley~\cite{wang2011analysis} analyze fly fishing rod casting dynamics. We focus on fixed-end cable manipulation and do not attempt to model cable physics, and we aim to control cables dynamically via vaulting, knocking, and weaving, and cannot rely on the same dragging motion primitives as in Kim and Shell~\cite{kim2016using}.

We propose an approach that is related to TossingBot~\cite{zeng_tossing_2019} where a robot learns to toss objects to target bins. TossingBot uses RGBD images and employs fully convolutional neural networks~\cite{fcn_2015} to learn dense, per-pixel values of tossing actions, and learns a residual velocity of tossing on top of a physics model in order to jointly train a grasping and tossing policy through trial and error. We take a similar approach in dynamically manipulating by reducing the learning process to learn only the joint angles of the apex of the motion. We do not rely on a physics model of the dynamic manipulation motion, and instead, leverage a quadratic programming-based motion-planning primitive to generate a fast ballistic motion that minimizes the jerk while moving through the learned apex point.
